\documentclass[11pt,a4paper]{article}
\usepackage[T1]{fontenc}
\usepackage{lmodern}
\usepackage[margin=1.08in]{geometry}
\usepackage{amsmath,amssymb,amsthm,mathtools}
\usepackage{microtype}
\usepackage[hidelinks]{hyperref}
\newtheorem{theorem}{Theorem}[section]
\newtheorem{lemma}[theorem]{Lemma}
\newcommand{\R}{\mathbb R}
\newcommand{\Om}{\Omega}
\newcommand{\dd}{\,\mathrm d}
\title{A positive equilibrium at the critical Dirichlet eigenvalue}
\author{Kaiwen Jin\\[2pt]\small\texttt{yui.ui.siki@gmail.com}}
\date{October 15, 2026}
\begin{document}
\maketitle
\begin{abstract}
We construct a connected unbounded uniformly smooth domain in
$\R^{21}$ whose bottom Dirichlet eigenvalue is $1$, but on which
$-\Delta u=u-u^2$ has a bounded positive solution with zero boundary values.
Thus equality between the spectral threshold and the derivative of a
strict KPP reaction at zero does not force extinction on all uniformly
smooth domains. The domain is a slight widening of a slab of width $\pi$.
Horizontal Hardy energy controls the transverse spectral deficit, while
an explicit positive subsolution survives the quadratic reaction loss.
\end{abstract}

\section{The critical existence question}
Consider the Dirichlet equilibrium equation
\begin{equation}\label{eq:problem}
 -\Delta u=f(u)\quad\text{in }\Om,\qquad u=0\quad\text{on }\partial\Om,
\end{equation}
where $\Om\subset\R^d$ is connected and unbounded. We use the spectral
normalization
\begin{equation}\label{eq:rayleigh}
 \lambda_1(\Om)=\inf_{0\ne\phi\in C_c^\infty(\Om)}
 \frac{\int_\Om|\nabla\phi|^2}{\int_\Om\phi^2}.
\end{equation}
For a reaction satisfying $f(0)=f(1)=0$, $f'(0)>0$, positivity on $(0,1)$,
negativity above $1$, and strict decrease of $f(s)/s$ on $(0,1]$, the
critical case is $\lambda_1(\Om)=f'(0)$. In the population interpretation,
this balances linear growth against diffusive loss.

Berestycki and Graham prove existence when $\lambda_1(\Om)<f'(0)$ and
nonexistence when $\lambda_1(\Om)>f'(0)$, as well as critical nonexistence
on bounded or periodic domains \cite[Theorem~1.5]{BG}. Their
Conjecture~1.7 asks for critical nonexistence in general domains. Uniform
boundary regularity is part of that setting: the boundary graph charts
have a common radius and uniformly bounded $C^{2,\gamma}$ norms.

Rossi and Tough establish dimension-dependent critical behavior under
quantitative conditions on the nonlinear loss \cite{RT}. Their
inverse-square potential construction involves the Hardy constant
\cite[Theorem~4.1]{RT}. They also obtain nonuniqueness for the homogeneous
logistic equation on locally, but not uniformly, analytic domains
\cite[Remark~4.2(e) and Section~7]{RT}. We give an explicit critical
existence example with uniform boundary regularity and a homogeneous reaction.

Write $(x,z)\in\R^{20}\times\R$, and define
\begin{equation}\label{eq:domain}
 t(x)=1000+|x|^2,\qquad b(x)=\frac{50}{t(x)},\qquad
 h(x)=\arccos(-b(x)),\qquad
 \Om=\{(x,z):|z|<h(x)\}.
\end{equation}
To make the global regularity of the reaction explicit, take
\begin{equation}\label{eq:reaction}
 f(s)=\begin{cases}
 s-s^2,&0\le s\le1,\\
 s^{-1}-1,&s\ge1.
 \end{cases}
\end{equation}
The two branches agree in value and first derivative at $1$; the function
and its derivative are bounded, and its derivative is globally Lipschitz.
It satisfies all the reaction conditions above.

\begin{theorem}\label{thm:main}
For every fixed $0<\gamma\le1$, the domain $\Om$ in \eqref{eq:domain} is
connected, unbounded and uniformly $C^{2,\gamma}$, and the reaction
\eqref{eq:reaction} belongs to $C^{1,\gamma}([0,\infty))$. Moreover,
\[
 \lambda_1(\Om)=f'(0)=1,
\]
and \eqref{eq:problem} admits a solution
$U\in C^2(\Om)\cap C(\overline\Om)$ with $0<U\le1$.
\end{theorem}

Here $\lambda_1$ is precisely the compact-test infimum
\eqref{eq:rayleigh}, and the boundary condition holds at every point of
$\partial\Om$. Since $U\le1$, it solves the classical logistic equation.
This single example answers the general critical nonexistence question in
the negative. The dimension and constants are convenient choices; no optimal
dimension is asserted.

The domain contains the critical slab of width $\pi$, but is wider by an
amount of order $|x|^{-2}$. The horizontal Hardy inequality pays for the
resulting decrease of the vertical eigenvalue. A cosine with a small
positive offset then supplies a nonlinear subsolution. Section~2 proves
the spectral equality, and Section~3 constructs the equilibrium and recovers
its boundary values. All functions and test spaces are real.

\section{Geometry and spectrum}
\begin{lemma}\label{lem:spectrum}
The domain in \eqref{eq:domain} is connected, unbounded and uniformly
$C^{2,\gamma}$ for $0<\gamma\le1$, and $\lambda_1(\Om)=1$.
\end{lemma}
\begin{proof}
The function $b$ has bounded derivatives of every order and takes values in
$(0,1/20]$. The same is true of the positive-order derivatives of $h$.
Also $h>\pi/2$ and
\[
 |\nabla h|\le2|\nabla b|\le\frac{200}{t^{3/2}}<\frac1{10}.
\]
Thus the two global graphs $z=\pm h(x)$ are uniformly separated and have
uniformly bounded $C^{2,\gamma}$ norms on fixed-radius charts. For example,
charts of horizontal radius $1/4$ have height variation less than $1/40$;
the other graph is more than $3$ away vertically. The midplane gives
connectedness, and $\Om$ contains
$S=\R^{20}\times(-\pi/2,\pi/2)$.

We first bound the transverse spectral deficit. For $0\le b\le1/20$, set
$\theta=2\arcsin(b)/\pi$. Since
$\arcsin(b)\le6b/5$ and $\pi>3$,
\begin{equation}\label{eq:deficit}
 1-\left(\frac{\pi}{2\arccos(-b)}\right)^2
 =1-(1+\theta)^{-2}\le2\theta\le\frac85 b.
\end{equation}
The derivative bound $2/\sqrt3<6/5$ on $[0,1/2]$ proves the arcsine
estimate used here.

For $\psi\in C_c^\infty(\R^{20})$, expanding a square gives
\begin{equation}\label{eq:hardy}
 0\le\int_{\R^{20}}\left|\nabla\psi+9\frac{x}{|x|^2}\psi\right|^2
 =\int_{\R^{20}}|\nabla\psi|^2
   -81\int_{\R^{20}}\frac{\psi^2}{|x|^2}.
\end{equation}
Indeed, $\operatorname{div}(x/|x|^2)=18/|x|^2$. Integrating first outside
a ball of radius $\varepsilon$ justifies the computation at the origin:
the extra boundary term is $O(\varepsilon^{18})$ and tends to zero.

Let $\phi\in C_c^\infty(\Om)$ and extend it by zero to $\R^{21}$. Its
extension remains smooth. Apply \eqref{eq:hardy} to each horizontal slice,
and the interval Dirichlet inequality to each vertical slice. For clarity,
the latter follows by writing $p(z)=\cos(\pi z/(2H))$ on $(-H,H)$ and using
\[
 \int_{-H}^H\left(|g'|^2-\frac{\pi^2}{4H^2}g^2\right)
 =\int_{-H}^H p^2\left|(g/p)'\right|^2\ge0
 \qquad(g\in C_c^\infty(-H,H)).
\]
With $r=|x|$, Fubini and \eqref{eq:deficit} therefore give
\begin{equation}\label{eq:lower-spectrum}
 \int_\Om|\nabla\phi|^2
 \ge\int_\Om\left(1-\frac{80}{1000+r^2}+\frac{81}{r^2}\right)\phi^2
 \ge\int_\Om\phi^2.
\end{equation}
The weighted term is integrable near $r=0$ in twenty dimensions. This
proves $\lambda_1(\Om)\ge1$.

For the reverse bound, choose $p_\varepsilon\in C_c^\infty(-\pi/2,\pi/2)$
converging to $\cos z$ in $H^1$. Such functions are obtained by cutting off
within distance $2\varepsilon$ of the endpoints; there $\cos z=O(\varepsilon)$,
so the derivative error tends to zero in $L^2$. For any nonzero
$\eta\in C_c^\infty(\R^{20})$, the product
$\eta(x/L)p_\varepsilon(z)$ is a compact test in $S\subset\Om$ and has
Rayleigh quotient
\[
 L^{-2}\frac{\int|\nabla\eta|^2}{\int\eta^2}
 +\frac{\int|p_\varepsilon'|^2}{\int p_\varepsilon^2}.
\]
Letting $L\to\infty$ and then $\varepsilon\to0$ proves the upper bound.
\end{proof}

\section{A positive equilibrium}
\begin{lemma}\label{lem:subsolution}
The smooth function
\begin{equation}\label{eq:subsolution}
 v(x,z)=\frac{\cos z}{t(x)}+\frac{50}{t(x)^2}
\end{equation}
is positive in $\Om$, vanishes on $\partial\Om$, is bounded by $21/20000$,
and satisfies $-\Delta v<v-v^2$ in $\Om$.
\end{lemma}
\begin{proof}
Positivity and the boundary values follow from
$\cos h(x)=-50/t(x)$ and $h<\pi$. With $r=|x|$ and $t=1000+r^2$,
direct differentiation yields
\[
 \Delta_x(t^{-1})=-40t^{-2}+8r^2t^{-3},\qquad
 \Delta_x(50t^{-2})=-4000t^{-3}+1200r^2t^{-4}.
\]
The first coefficient is negative and at least $-40t^{-2}$.
The vertical cosine term cancels in $\Delta v+v$, so
\[
 \Delta v+v
 =\Delta_x(t^{-1})\cos z+\Delta_x(50t^{-2})+50t^{-2}
 \ge\frac{10}{t^2}-\frac{4000}{t^3}\ge\frac6{t^2}.
\]
On the other hand, $0<v\le(21/20)t^{-1}$ and
$v^2\le441/(400t^2)<6/t^2$.
\end{proof}

\begin{proof}[Proof of Theorem~\ref{thm:main}]
The reaction properties follow from \eqref{eq:reaction}, and
Lemma~\ref{lem:spectrum} proves the geometric and spectral claims.
It remains to construct a classical Dirichlet solution above the function
$v$ of Lemma~\ref{lem:subsolution}.

Fix $R>1$ and put $D_R=\Om\cap\{|x|<R\}$. This is a bounded Lipschitz
domain: the map $(x,y)\mapsto(x,h(x)y)$ sends
$B_R^{20}\times(-1,1)$ smoothly onto it, including its transverse corners.
Set $w_0=v$ and define $w_{k+1}$ by
\begin{equation}\label{eq:iteration}
 (-\Delta+1)w_{k+1}=G(w_k)\quad\text{in }D_R,\qquad
 w_{k+1}-v\in H^1_0(D_R),\qquad G(s)=2s-s^2.
\end{equation}
All boundary data here are weak traces. After subtracting $v$, the bilinear
form $\int_{D_R}(\nabla a\cdot\nabla c+ac)$ is the full $H^1_0$ inner
product. Lax--Milgram therefore defines every iterate whose right side is
in $L^2$ \cite[Theorem~4.20]{HunterPDE}. Weak comparison is obtained by
testing a difference with its positive part, which has zero trace when
the boundary difference is nonpositive; the Sobolev and trace facts used
here hold on bounded Lipschitz domains \cite[Section~1.1]{FRRO}.
Since $G$ is nondecreasing on $[0,1]$, $v$ is a subsolution, and $1$ is a
supersolution, induction gives
\begin{equation}\label{eq:monotone}
 v\le w_k\le w_{k+1}\le1.
\end{equation}

The limit in this iteration preserves the nonlinear equation. If $w(F)$
denotes the affine Dirichlet solution with right side $F$, testing the
difference of two equations gives
\[
 \|w(F)-w(F')\|_{H^1(D_R)}\le\|F-F'\|_{L^2(D_R)}.
\]
The monotone limit $U_R$ in \eqref{eq:monotone} satisfies
$G(w_k)\to G(U_R)$ in $L^2(D_R)$ by dominated convergence. The displayed
stability estimate then identifies $U_R$ with $w(G(U_R))$ in $H^1$.
Thus $U_R-v\in H^1_0(D_R)$ and
$-\Delta U_R=U_R-U_R^2$ weakly.

On each fixed interior ball the right side is bounded by $1/4$ in absolute
value. Interior regularity first gives uniform $C^{1,1/2}$ bounds and then,
by polynomial composition and Schauder estimates, uniform $C^{2,1/2}$
bounds on smaller balls \cite[Proposition~2.18 and Corollary~2.16]{FRRO}.
Taking $R\to\infty$ along a diagonal subsequence gives a limit
$U\in C^2(\Om)$ solving the equation, with $v\le U\le1$.

To retain the zero boundary values, write
$q(x,z)=\cos z+b(x)$ and work in the two collars $0<q<1/4$. There define
\[
 W(x,z)=\frac{1-e^{-4q(x,z)}}{1-e^{-1}}.
\]
Since $|\cos z|=|q-b|\le1/4$ and $\Delta_x b<0$,
\[
 |\nabla q|^2\ge\frac{15}{16},\qquad \Delta q\le1,
 \qquad
 -\Delta W
 =\frac{4e^{-4q}}{1-e^{-1}}
   \left(4|\nabla q|^2-\Delta q\right)
 \ge\frac{11}{e-1}>\frac14.
\]
The function of $q$ defining $W$ is concave, equals $0$ at $q=0$ and $1$
at $q=1/4$, and hence $W\ge4q$ throughout the collar. On the boundary of
$D_R\cap\{0<q<1/4\}$ the traces satisfy $U_R\le W$: both are zero at
$q=0$, $U_R\le1=W$ at $q=1/4$, and the artificial side $r=R$ has trace
$v=q/t\le W$. Each collar component is bounded Lipschitz, so weak
comparison with $-\Delta U_R\le1/4<-\Delta W$ gives $U_R\le W$ there.
This bound is independent of $R$. Passing to the limit yields
$0<U\le W$ near the original boundary. As $q\to0$, also $W\to0$,
which proves that $U$ extends continuously by zero to every point of
$\partial\Om$. Finally, $U\le1$ means $f(U)=U-U^2$, as required.
\end{proof}

\paragraph{AI use.}
OpenAI Codex was used to develop and check the construction and to prepare
the exposition. Independent mathematical review is requested.

\begin{thebibliography}{9}
\bibitem{BG}
H.~Berestycki and C.~Graham,
\emph{The steady states of strong-KPP reactions in general domains},
arXiv:2212.06611v1 (2022),
\url{https://arxiv.org/abs/2212.06611v1}.

\bibitem{RT}
L.~Rossi and O.~Tough,
\emph{Sharp dimension threshold for unique FKPP steady states and BBM recurrence},
arXiv:2610.17121v1 (2026),
\url{https://arxiv.org/abs/2610.17121v1}.

\bibitem{FRRO}
X.~Fern\'andez-Real and X.~Ros-Oton,
\emph{Regularity Theory for Elliptic PDE},
arXiv:2301.01564v1 (2023),
\url{https://arxiv.org/abs/2301.01564v1}.

\bibitem{HunterPDE}
J.~K.~Hunter,
\emph{Notes on Partial Differential Equations},
University of California, Davis, lecture notes,
\url{https://math.ucdavis.edu/~hunter/pdes/pde_notes.pdf},
accessed October 15, 2026.
\end{thebibliography}
\end{document}
